\section{Characters and phase invariance}\label{sec:preliminaries}

The Hermitian inner product is conjugate-linear in its first argument.
We index the coordinates of \(\C^6\) by \(0,\ldots,5\), and write
\(e_0,\ldots,e_5\) for the standard coordinate vectors.
For a square Hadamard matrix, $H^*H=6I$ also implies $HH^*=6I$.
The equivalence relation and normalized character $g_H$ are those of
Section~\ref{sec:introduction}. Every Hadamard matrix can be
\emph{dephased}, meaning that its first row and first column are all ones,
by multiplying rows and columns by elements of $\T$.

A \emph{charge} is a vector \(a\in\Z^6\) with \(\sum_i a_i=0\).
Its order and character are
\[
 \ord(a)=\sum_{a_i>0}a_i,
 \qquad x^a=\prod_{i=0}^5x_i^{a_i}\quad(x\in\T^6).
\]
For the column characters in \eqref{eq:g-definition}, put
\[
 D=\{e_i-e_j:0\le i,j\le5,\ i\ne j\}.
\]
We retain $\alpha=(1,1,1,-1,-1,-1)$ from the introduction.
In particular, \(|g_H(a)|\le1\), \(g_H(0)=1\), and
\(g_H(-a)=\overline{g_H(a)}\). A permutation \(p\) acts by
\((P_pa)_i=a_{p(i)}\). When we write a permuted charge, all coordinates are
permuted together; we also abbreviate \(P_pa\) as \(pa\).

Balanced charges make \(g_H\) invariant under column phases. Multiplying
row \(i\) by \(u_i\) multiplies \(g_H(a)\) by \(u^a\). Thus products
whose total charge is zero are unchanged by row phases as well.
These elementary observations will justify dephasing and the invariant
moments below.

